\subsection{Gate-length and top-silicon-thickness study: completed transfer sweeps}
Two numerical student-ID digits, $A=2$ and $B=2$, give $\delta=(A+B)/2=2$. Hence
$L_G=52,82,102,252,502,1002$ nm. The upper-silicon thicknesses were
$T_{\mathrm{SI}}=30$ and 10 nm. All remaining SProcess inputs used the
nominal baseline values. Twelve distinct SProcess geometries and 24
SDevice $I_D$--$V_G$ sweeps (at $V_{DS}=0.1$ and 1.2 V) were completed
in Sentaurus Workbench J-2014.09. The simulation job statuses and
the individual $I_D$--$V_G$ datasets were checked.
The original 24 export files contain 1,310 total rows. After retaining
the last point of each repeated $V_{GS}=-0.5$ V drain-initialization
ramp, the analyzed data consist of 35 gate voltages per transfer
curve ($-0.5$ through 1.2 V in 0.05-V steps), i.e., 840 points.

\paragraph{Extraction method and limits.}
Electrical metrics were calculated externally using Python (pandas
and NumPy) directly from the Sentaurus Visual exports.
The reference-current gate voltage $V_{\mathrm{TH,CC}}$ is the crossing
of $I_D=10^{-7}$ in the \emph{exported current units}, found by
linear interpolation in $\log_{10} I_D$ between adjacent positive
current values. Its value depends on the selected, not
width-normalized current criterion. DIBL is calculated from the
matched criterion at both drain biases,
\[
\mathrm{DIBL}_{\mathrm{CC}}=
1000\,\frac{V_{\mathrm{TH,CC}}(0.1\mathrm{\,V})-
V_{\mathrm{TH,CC}}(1.2\mathrm{\,V})}{1.1\mathrm{\,V}}
\quad(\mathrm{mV/V}).
\]
The reported SS is $1000$ divided by the fitted slope of
$\log_{10}I_D$ against $V_{GS}$ over the same six gate-bias points,
$-0.20$ through $+0.05$ V, at $V_{DS}=0.1$ V.
For severely short-channel devices this is an \emph{apparent
subthreshold slope} of a fixed window; it is not proof that a true
weak-inversion region has been captured. In particular, at
$T_{\mathrm{SI}}=30$ nm and $L_G=52$ nm, the high-drain-bias current
is already $3.348\times10^{-5}$ at the minimum swept gate voltage
$-0.5$ V. Consequently $I_D=10^{-7}$ is never reached in that
high-drain sweep; its $V_{\mathrm{TH,CC}}(1.2)$ and DIBL are
\textbf{not reported}, not extrapolated.

Here $I_{\mathrm{ON}}$ is defined at $V_{GS}=V_{DS}=1.2$ V and
$I_{\mathrm{OFF}}$ at $V_{GS}=0$, $V_{DS}=1.2$ V.
Current values remain in the simulator-exported units until
the two-dimensional current-width normalization is checked.
Both quantities are extracted directly from the simulated
bias points.

\paragraph{Quantitative metric comparison.}
\begin{table}[H]\centering\scriptsize
\caption{Computed gate-length sweep metrics. \texttt{N/A} means the
constant-current threshold was outside the simulated gate-bias
range. Currents are in exported Sentaurus units.}
\resizebox{\linewidth}{!}{%
\begin{tabular}{@{}rrrrrrrr@{}}\toprule
$T_{\mathrm{SI}}$ (nm) & $L_G$ (nm) & $V_{\mathrm{TH,0.1}}$ (V) &
$V_{\mathrm{TH,1.2}}$ (V) & SS (mV/dec) &
$I_{\mathrm{ON}}$ & $I_{\mathrm{OFF}}$ & DIBL (mV/V)\\\midrule
30&52&$-0.4037$&N/A&269.3&$1.3543\times10^{-3}$&$1.0805\times10^{-4}$&N/A\\
30&82&0.0158&$-0.3084$&113.7&$1.0471\times10^{-3}$&$7.3799\times10^{-6}$&294.7\\
30&102&0.0843&$-0.0871$&90.1&$9.3697\times10^{-4}$&$6.0713\times10^{-7}$&155.8\\
30&252&0.1778&0.1599&63.6&$5.7877\times10^{-4}$&$3.3930\times10^{-10}$&16.2\\
30&502&0.2063&0.1991&61.5&$3.7524\times10^{-4}$&$6.1314\times10^{-11}$&6.6\\
30&1002&0.2298&0.2263&60.9&$2.2549\times10^{-4}$&$2.1310\times10^{-11}$&3.2\\\midrule
10&52&0.0172&$-0.1692$&91.2&$1.2129\times10^{-3}$&$3.6018\times10^{-6}$&169.4\\
10&82&0.1092&0.0559&69.2&$9.9068\times10^{-4}$&$1.6965\times10^{-8}$&48.4\\
10&102&0.1288&0.0966&65.5&$8.9849\times10^{-4}$&$3.6369\times10^{-9}$&29.2\\
10&252&0.1726&0.1631&61.5&$5.7735\times10^{-4}$&$2.3113\times10^{-10}$&8.7\\
10&502&0.1958&0.1911&61.0&$3.8161\times10^{-4}$&$7.5608\times10^{-11}$&4.2\\
10&1002&0.2175&0.2149&60.8&$2.3271\times10^{-4}$&$3.0831\times10^{-11}$&2.4\\
\bottomrule\end{tabular}}
\end{table}

\paragraph{Five required comparisons.}
The following five plots use the numerical archives
\texttt{data/part2\_tsi30\_metrics.csv} and
\texttt{data/part2\_tsi10\_metrics.csv}. The raw data are plotted
without smoothing. Missing $V_{\mathrm{TH}}$ and DIBL values for
the 30-nm, 52-nm case are explicitly omitted.

% Common axis choice: logarithmically spaced display for the six tested LG values.
\begin{figure}[H]\centering
\begin{tikzpicture}
\begin{axis}[width=.95\linewidth,height=6.6cm,xmode=log,log basis x=10,
xtick={52,82,102,252,502,1002},xticklabels={52,82,102,252,502,1002},
xlabel={$L_G$ (nm)},ylabel={$V_{\mathrm{TH,CC}}$ (V)},
grid=both,legend style={font=\scriptsize},legend pos=south east,unbounded coords=jump]
\addplot+[mark=*,thick] table[x=LG_nm,y=VTH_0p1_V,col sep=comma] {data/part2_tsi30_metrics.csv};
\addlegendentry{30 nm, $V_{DS}=0.1$ V}
\addplot+[mark=square*,thick,dashed] table[x=LG_nm,y=VTH_1p2_V,col sep=comma] {data/part2_tsi30_metrics.csv};
\addlegendentry{30 nm, $V_{DS}=1.2$ V}
\addplot+[mark=triangle*,thick] table[x=LG_nm,y=VTH_0p1_V,col sep=comma] {data/part2_tsi10_metrics.csv};
\addlegendentry{10 nm, $V_{DS}=0.1$ V}
\addplot+[mark=diamond*,thick,dashed] table[x=LG_nm,y=VTH_1p2_V,col sep=comma] {data/part2_tsi10_metrics.csv};
\addlegendentry{10 nm, $V_{DS}=1.2$ V}
\end{axis}
\end{tikzpicture}
\caption{Constant-current $V_{\mathrm{TH}}$ against $L_G$, for both
silicon thicknesses and both drain biases.}
\end{figure}

\begin{figure}[H]\centering
\begin{tikzpicture}
\begin{axis}[width=.95\linewidth,height=6.4cm,xmode=log,log basis x=10,
xtick={52,82,102,252,502,1002},xticklabels={52,82,102,252,502,1002},
xlabel={$L_G$ (nm)},ylabel={SS (mV/dec)},grid=both,
legend pos=north east,unbounded coords=jump]
\addplot+[mark=*,thick] table[x=LG_nm,y=SS_0p1_mV_dec,col sep=comma] {data/part2_tsi30_metrics.csv};
\addlegendentry{$T_{\mathrm{SI}}=30$ nm}
\addplot+[mark=square*,thick] table[x=LG_nm,y=SS_0p1_mV_dec,col sep=comma] {data/part2_tsi10_metrics.csv};
\addlegendentry{$T_{\mathrm{SI}}=10$ nm}
\end{axis}
\end{tikzpicture}
\caption{Same-window $V_{DS}=0.1$ V apparent subthreshold swing versus gate length.}
\end{figure}

\begin{figure}[H]\centering
\begin{tikzpicture}
\begin{axis}[width=.95\linewidth,height=6.4cm,xmode=log,log basis x=10,
xtick={52,82,102,252,502,1002},xticklabels={52,82,102,252,502,1002},
xlabel={$L_G$ (nm)},ylabel={$I_{\mathrm{ON}}$ (exported units)},
grid=both,legend pos=north east,unbounded coords=jump]
\addplot+[mark=*,thick] table[x=LG_nm,y=ION_exported,col sep=comma] {data/part2_tsi30_metrics.csv};
\addlegendentry{$T_{\mathrm{SI}}=30$ nm}
\addplot+[mark=square*,thick] table[x=LG_nm,y=ION_exported,col sep=comma] {data/part2_tsi10_metrics.csv};
\addlegendentry{$T_{\mathrm{SI}}=10$ nm}
\end{axis}
\end{tikzpicture}
\caption{ON current at $V_{GS}=V_{DS}=1.2$ V.}
\end{figure}

\begin{figure}[H]\centering
\begin{tikzpicture}
\begin{semilogyaxis}[width=.95\linewidth,height=6.4cm,xmode=log,log basis x=10,
xtick={52,82,102,252,502,1002},xticklabels={52,82,102,252,502,1002},
xlabel={$L_G$ (nm)},ylabel={$I_{\mathrm{OFF}}$ (exported units)},
grid=both,legend pos=north east,ymin=1e-11,ymax=1e-3,unbounded coords=jump]
\addplot+[mark=*,thick] table[x=LG_nm,y=IOFF_exported,col sep=comma] {data/part2_tsi30_metrics.csv};
\addlegendentry{$T_{\mathrm{SI}}=30$ nm}
\addplot+[mark=square*,thick] table[x=LG_nm,y=IOFF_exported,col sep=comma] {data/part2_tsi10_metrics.csv};
\addlegendentry{$T_{\mathrm{SI}}=10$ nm}
\end{semilogyaxis}
\end{tikzpicture}
\caption{OFF current at $V_{GS}=0$ and $V_{DS}=1.2$ V.}
\end{figure}

\begin{figure}[H]\centering
\begin{tikzpicture}
\begin{axis}[width=.95\linewidth,height=6.4cm,xmode=log,log basis x=10,
xtick={52,82,102,252,502,1002},xticklabels={52,82,102,252,502,1002},
xlabel={$L_G$ (nm)},ylabel={DIBL (mV/V)},grid=both,
legend pos=north east,unbounded coords=jump]
\addplot+[mark=*,thick] table[x=LG_nm,y=DIBL_mV_V,col sep=comma] {data/part2_tsi30_metrics.csv};
\addlegendentry{$T_{\mathrm{SI}}=30$ nm}
\addplot+[mark=square*,thick] table[x=LG_nm,y=DIBL_mV_V,col sep=comma] {data/part2_tsi10_metrics.csv};
\addlegendentry{$T_{\mathrm{SI}}=10$ nm}
\end{axis}
\end{tikzpicture}
\caption{Matched-reference-current DIBL against gate length.
The missing 30-nm, 52-nm point was outside the gate sweep's threshold
crossing range.}
\end{figure}

\paragraph{Interpretation and the first significant departure.}
Both devices show near-ideal apparent SS at long gate lengths
($\simeq61$ mV/dec at 1002 nm). Define a significant SS departure
as a \emph{greater than 10\% increase} from the respective
1002-nm reference while reducing $L_G$. With this explicit criterion,
the first tested $L_G$ satisfying it is 102 nm for the 30-nm
film ($90.10$ versus $60.87$ mV/dec, a 48.0\% increase) and
82 nm for the 10-nm film ($69.16$ versus $60.84$ mV/dec, a
13.7\% increase). DIBL responds even earlier in \emph{relative}
terms, increasing from $3.17$ to $6.58$ mV/V between 1002
and 502 nm in the 30-nm film, and from $2.39$ to $4.21$
mV/V in the 10-nm film; the corresponding absolute changes
are still small at those longer lengths.

At 82 nm, switching from 30 to 10 nm top silicon reduces
$I_{\mathrm{OFF}}$ by a factor of approximately 435,
reduces DIBL from $294.7$ to $48.4$ mV/V,
and improves apparent SS from $113.7$ to $69.2$ mV/dec.
The trade-off is a 5.39\% decrease in $I_{\mathrm{ON}}$.
At 52 nm the ON-current reduction is 10.44\%, while OFF
current falls by a factor of approximately 30. This refutes
a claim of \emph{no} ON-current penalty. At longer
gate lengths the leakage advantage is small or reversed,
so the benefit is conditional on the geometry and bias.
This performance trend is consistent with improved
electrostatic control in the thinner SOI film.

\paragraph{Depletion interpretation and remaining Chapter 4 evidence.}
The observed SS and DIBL improvements are consistent with
stronger gate control and suppression of drain-induced
barrier lowering in thin-film SOI. However, transfer-curve
metrics \emph{alone} cannot prove that a film is fully
depleted versus partially depleted. The requested 2D
electrostatic-potential comparison at $V_{GS}=0$ V and
$V_{DS}=1.2$ V, supplemented by a channel-center charge
profile, is still needed for a spatially resolved claim.
Also still pending are the shortest-versus-longest
$L_G$ $I_D$--$V_D$ overlays for \emph{each} $T_{\mathrm{SI}}$.
These plots require further bias sweeps beyond the
completed transfer-characteristic batch.
